\section{Matter bridge: solitons as self-adjoint eigenproblems}
\label{sec:matter}

\subsection{Spherical material coordinates (coordinate change only)}
\label{subsec:spherical-material-coords}

Localized particle-like states are organized most naturally in a spherical chart on the
intrinsic brane domain. We record the relevant coordinate change here, before linearizing.

\paragraph{Important: this is only a coordinate change.}
All dynamical statements remain those of the coordinate-free formulation
$\vecX:\Omega\times\R\to\R^4$ with induced metric $g_{ij}=\partial_i\vecX\cdot\partial_j\vecX$.
We introduce spherical coordinates on the \emph{intrinsic/material} domain as a convenient chart for localized
excitations; this does \emph{not} constrain the embedding $\vecX$.

Let $(r,\theta,\varphi)$ be the standard spherical chart on (a subset of) $\R^3$, with Jacobian
$J(r,\theta,\varphi)=r^2\sin\theta$.
Writing the intrinsic coordinate as $q^a\in\{r,\theta,\varphi\}$ and denoting derivatives in this chart by $\partial_a$,
the equations of motion (\Cref{eq:eom} of Paper~I) become
\begin{equation}
  \rho_m\,\partial_{tt}X^A
  =\frac{1}{J(q)}\,\partial_a\!\Big(J(q)\,P_A^{\;a}\Big),
  \qquad A=1,\dots,4,
  \label{eq:eom-spherical}
\end{equation}
which is exactly the Cartesian divergence form (\Cref{eq:eom} of Paper~I) written with the standard Jacobian factor.

\subsection{Linearization about symmetric backgrounds}
\label{subsec:linearization}

Let $\vecX_0(x)$ be a time-independent background embedding representing a localized deformation
region, taken to be spherically symmetric in the intrinsic brane coordinates
($r=|x|$ with $r\gg a$ for the continuum limit to apply).
Linearize the equations of motion (\Cref{eq:eom} of Paper~I) about $\vecX_0$ by writing
$\vecX(x,t)=\vecX_0(x)+\bm{\eta}(x,t)$ with small perturbation $\bm{\eta}$.
Because the dynamics derives from the action (\Cref{eq:action} of Paper~I), the linearized
operator is the second variation of $S$ and is therefore self-adjoint with respect to the
$\rho_m$-weighted $L^2$ inner product on $\Omega$; orthogonality and completeness of eigenmodes
follow as standard consequences.

\subsection{Angular structure: vector spherical harmonics}
\label{subsec:vsh}

The substrate is the 6-neighbor cubic lattice, whose point group is $O_h$ rather than $SO(3)$.
At scales $\gg a$, however, long-wavelength continuum modes recover approximate isotropy with
residual cubic corrections of order 7.5\% in longitudinal speed at $\alpha=0.2$.
Solitons in this theory are expected to live in the geometric-quartic-stabilized regime with
$\ell_{\mathrm{soliton}}\gg a$, so the emergent approximate $SO(3)$ is the right symmetry to
organize their angular structure, with $O_h$ corrections entering as small splittings of $SO(3)$
multiplets.

For a spherically symmetric background the linearized operator commutes with the action of $SO(3)$
on the intrinsic coordinates at leading order in $a/\ell_{\mathrm{soliton}}$.
For a scalar field this gives the familiar separation into radial and angular problems with
eigenfunctions $Y_{\ell m}(\theta,\varphi)$.
The lateral triplet $\xi^i\in\R^3$ and its complex-envelope promotion $\Psi\in\C^3$ from
Paper~III carry an internal vector index $i$ in addition to spatial dependence.
The correct angular basis is therefore the \emph{vector spherical harmonic} (VSH) basis
$\mathbf{Y}^i_{JLM}$, where $J$ is total angular momentum, $L$ is orbital angular momentum,
and $L\in\{J-1,J,J+1\}$:
\begin{equation}
  \xi^i(r,\theta,\varphi)
  \;=\; \sum_{J,L,M} f_{JLM}(r)\,\mathbf{Y}^i_{JLM}(\theta,\varphi).
  \label{eq:vsh-decomposition}
\end{equation}
The choice of $(J,L)$ specifies how the internal index $i$ is glued to the spatial angular
coordinates --- it is the choice of how the triplet's $U(3)$ gauge structure (Paper~III) couples to
the geometry of the localized mode.
Configurations with $J=0$ are spherically symmetric only under the \emph{combined} rotation that
acts on both spatial and internal indices together; this combined rotation is the physically
meaningful notion of spherical symmetry for a triplet field.

\subsection{Hedgehog ansatz and Skyrme-twisted extension}
\label{subsec:hedgehog}

The canonical baryon ansatz is the \emph{hedgehog}
\begin{equation}
  \xi^i(\mathbf{x}) \;=\; f(r)\,\hat x^i,
  \label{eq:hedgehog}
\end{equation}
with $J=0$, $L=1$: the internal index $i$ is locked to the spatial direction $\hat x^i$, so the
configuration is invariant only under the combined diagonal $SO(3)$.
The entire nontrivial structure of the hedgehog lives in the $SU(3)$ traceless sector of Paper~III:
the $U(1)$ trace $f(r)\,(\hat x^1+\hat x^2+\hat x^3)$ averages to zero on the sphere, giving a
trace-neutral far field (neutron-like far-field holonomy).

Topological stabilization at winding number $B=1$ is achieved by extending the hedgehog into a
\emph{Skyrme-twisted} configuration in which the embedding-amplitude direction $X^4$ also rotates
radially:
\begin{equation}
  \bigl(\xi^i,u\bigr)(\mathbf{x})
  \;=\; \bigl(f(r)\sin F(r)\,\hat x^i,\;f(r)\cos F(r)\bigr),
  \label{eq:skyrme-twist}
\end{equation}
with the profile $F:[0,\infty)\to[\pi,0]$ enforcing the boundary conditions
$F(0)=\pi$ and $F(\infty)=0$.
This boundary condition gives $B=1$ topological winding in $\pi_3(SU(3))=\mathbb{Z}$, consistent
with Paper~III \Cref{subsec:topology}.
Adding a nonzero $L=0$ scalar admixture shifts the $U(1)$ far-field holonomy from trace-neutral
to proton-like.

\subsection{Soliton labels}
\label{subsec:soliton-labels}

A localized 3D mode is fully classified at the soliton layer by the labels
\begin{equation}
  \bigl(J,\,P;\;\mathrm{SU}(3)\text{-irrep};\;Q_{U(1)};\;B_{\mathrm{winding}}\bigr),
  \label{eq:soliton-labels}
\end{equation}
where:
\begin{itemize}
  \item $J$ is total angular momentum from the emergent approximate $SO(3)$ at scales $\gg a$;
  it refines to an $O_h$ irrep at sub-leading order in $a/\ell_{\mathrm{soliton}}$.
  \item $P$ is parity, the discrete element of the emergent rotation group.
  \item The $SU(3)$ irrep and $U(1)$ charge come from the $U(3)$ decomposition (Paper~III
  \Cref{subsec:u3-decomp}).
  \item $B_{\mathrm{winding}}$ is the topological winding number in $\pi_3(SU(3))=\mathbb{Z}$.
\end{itemize}
These labels are not independent: the VSH coupling sets a Clebsch--Gordan relationship between
$SO(3)$ and $U(3)$ representations, and the charge-triality locking $q\equiv -t\pmod{3}$
(Paper~III \Cref{subsec:u3-decomp}) constrains $Q_{U(1)}$ given the $SU(3)$ irrep.

\subsection{Stability: Derrick's theorem on the discrete substrate}
\label{subsec:derrick}

In 3D, a standard Derrick scaling argument \citep{Derrick1964} rules out stable finitely-extended
solitons in a purely quadratic (linear) theory: under $x\mapsto\lambda x$, the gradient energy
$E_2$ scales as $\lambda^{1}$ (favoring expansion for $\lambda>1$) while a quartic energy $E_4$
scales as $\lambda^{-1}$ (resisting expansion).
Without $E_4>0$, the soliton expands without bound.
The present substrate breaks the standard Derrick counting in two complementary ways.

\paragraph{Anti-collapse: lattice UV cutoff.}
Derrick's argument assumes continuous rescaling $\lambda\in(0,\infty)$.
On a discrete lattice with spacing $a$, however, $\lambda$ is bounded below: a soliton cannot
shrink past one cell.
A configuration with energy fully concentrated at a single node is not a local energy minimum ---
the spring network immediately disperses it outward.
The lattice spacing $a$ is the physical UV cutoff; no separate quartic repulsion is needed to block
collapse.
This is a genuinely discrete-substrate effect that has no direct continuum analogue.

\paragraph{Anti-expansion: geometric quartic.}
Under Derrick rescaling in 3D, the geometric quartic $W_4\propto k_s\alpha/a$ (Paper~I
\Cref{eq:W4}) scales as $\lambda^{-1}$.
Balancing $\lambda E_2$ against $\lambda^{-1}E_4$ gives equilibrium at
\begin{equation}
  \lambda^* = \sqrt{E_4/E_2},
  \label{eq:derrick-balance}
\end{equation}
translating to an equilibrium soliton half-radius
\begin{equation}
  \frac{R_h}{a}
  \;\approx\;
  \kappa\,\frac{A}{a}\,\sqrt{\frac{\alpha}{1-\alpha}},
  \label{eq:derrick-radius}
\end{equation}
where $\kappa$ is a numerical factor of order unity fixed by the Derrick balance in the specific
ansatz and $A$ is the peak amplitude.
At $\alpha=0.2$ and $A/a\approx 10$ this gives $R_h/a\approx 5$; the soliton sweet spot
$\alpha\approx 0.5$--$0.8$ gives $R_h/a\approx 8$--$14$ for the same amplitude, well above the
lattice scale.

Note that the quartic coefficient scales as $\propto\alpha/a$, not $\propto(1-\alpha)$ as a naive
StVK identification would give.
The correct sign and scaling are derived from the spring's geometric quartic in Paper~I
(\Cref{eq:W4}), which inverts the na\"\i ve StVK prediction.
Quantitative soliton sizing must use the spring law, not StVK.

\subsection{Periodic boundary conditions: mandatory requirement}
\label{subsec:periodic-bc}

Periodic boundary conditions (BC) are mandatory for all nonlinear soliton simulations.
The prestressed lattice (rest length $\alpha a < a$) is under tension; open/free BC allows the box
to contract from $Na$ to $N\alpha a$, which dominates all dynamics and collapses even the $A=0$
vacuum.
Periodic BC fixes the box at $N\cdot a$, holds the tension, and ensures the vacuum is stable
($A=0\Rightarrow E_{\mathrm{excess}}\equiv 0$).
All retracted 2026-06-05 runs used open BC; all future nonlinear runs must use periodic-clamped BC.

\subsection{Matter bridge open problems}
\label{subsec:matter-open}

\begin{enumerate}
  \item \textbf{Converged baryon eigenstate.}
  Obtain a converged, stable, localized $\pi_3$-winding solution using the periodic-clamped
  tensioned vacuum and an eigenstate approach (not seed-and-watch evolution).
  The hedgehog/Skyrme-twisted ansatz \eqref{eq:skyrme-twist} with $\alpha\approx 0.7$,
  $A/a\approx 10$ is the primary candidate.

  \item \textbf{Rigorous Derrick stability proof.}
  State Derrick's theorem precisely \citep{Derrick1964,MantonSutcliffe2004} for the present
  action, give the lattice-UV anti-collapse argument as a theorem (not just intuition), and tie
  the anti-expansion balance \eqref{eq:derrick-balance} to the effective field manifold of
  Paper~III.

  \item \textbf{Converged $U(1)$ vortex eigenstate.}
  Resolve the semilocal drift observed in the preliminary JFNK run (Paper~III
  \Cref{subsec:numerical-targets}) with branch-selection machinery for the semilocal vortex.

  \item \textbf{Two-soliton interaction potential.}
  Once individual stable eigenstates exist, measure the two-soliton interaction potential as a
  function of separation and polarization alignment, to test whether the geometric self-guidance
  mechanism produces an attractive or repulsive channel.

  \item \textbf{Spin-$\tfrac12$ from soliton holonomy.}
  The spinorial $4\pi$ aspect identified kinematically in Paper~III
  (\Cref{subsec:spinorial-holonomy}) must be realized by an actual localized mode: show that a
  converged soliton carries a $\mathbb{Z}_2$ holonomy ($\pi_1(SO(3))=\mathbb{Z}_2$) under a
  $2\pi$ spatial rotation, so that spin-$\tfrac12$ emerges at the soliton level rather than from
  the linear envelope.
\end{enumerate}